\documentclass[11pt]{article}
\usepackage[margin=1in]{geometry}
\usepackage{amsmath,amssymb,amsthm}
\usepackage[T1]{fontenc}
\usepackage{lmodern}
\usepackage{graphicx}
\usepackage[hidelinks]{hyperref}
\usepackage{xcolor}

\newtheorem{theorem}{Theorem}
\newtheorem{lemma}[theorem]{Lemma}
\newtheorem{corollary}[theorem]{Corollary}
\newcommand{\R}{\mathbb R}
\newcommand{\cH}{\mathcal H}
\newcommand{\SOS}{\operatorname{SOS}}
\newcommand{\E}{\mathbb E}

\title{Eventual Sums of Squares under Heat Flow:\\
The Coercive-Leading-Form Case}
\author{Candidate partial result for arXiv:2211.04416}
\date{13 August 2026}

\begin{document}
\maketitle

\begin{center}
\fcolorbox{black}{gray!8}{\parbox{0.88\linewidth}{
\textbf{Status: candidate partial result, likely valid pending human review.}
The source problem remains open for degenerate leading forms.}}
\end{center}

\section{Source problem}

Curto, di Dio, Korda, and Magron study the action of the polynomial heat
semigroup on nonnegative polynomials \cite{CDKM}.  On PDF page 26 they ask:

\begin{quote}
\textbf{Open Problem 5.1.}
Let $f\in\operatorname{Pos}(\R^2)\setminus\SOS(\R^2)$.  Is it true that
there is always a $T=T(f)>0$ such that
$e^{t\Delta}f\in\SOS(\R^2)$ for all $t\geq T$?
\end{quote}

\begin{figure}[h]
  \centering
  \includegraphics[width=\textwidth]{figures/open_problem_crop.png}
  \caption{Open Problem 5.1 and its context in arXiv:2211.04416, PDF page 26.}
\end{figure}

For $m,d\geq1$, let $\cH_{m,d}$ be the real vector space of degree-$d$
forms in $m$ variables, and let $\Sigma_{m,2d}$ be the cone of sums of
squares of forms in $\cH_{m,d}$.  A form is a \emph{strict sum of squares}
if it lies in the interior of $\Sigma_{m,2d}$; equivalently, it has a
positive-definite Gram matrix in a monomial basis of $\cH_{m,d}$.

\section{Proof intuition}

Parabolic scaling separates the highest homogeneous part $p$ of $f$ from
all lower-order terms.  After homogenization in a new variable $z$, the
normalized late-time heat orbit converges to $e^{z^2\Delta}p$.

If $p$ is a strict SOS, choose a Gram factorization
$p=\sum_jq_j^2$ in which the $q_j$ form a basis of all degree-$d$ forms.
The Gaussian formula for heat flow writes the limit as an average of the
squares of the translated forms
$q_j(x+\sqrt2zG)$.  Those translations span every degree-$d$ form in the
old variables together with $z$.  Consequently the averaged Gram matrix is
positive definite: the limiting form is an \emph{interior} SOS point, not
merely an SOS point.  Small lower-degree perturbations are therefore absorbed
for all sufficiently large times.

For two variables, every nonnegative binary form is SOS by Hilbert's binary
form theorem \cite{Hilbert}.  A positive-definite binary form is an interior
point of that cone, so the argument covers exactly the source problem's
coercive-leading-form case.  The remaining difficulty is the boundary case
where the leading form vanishes in a real projective direction.

\section{A Gaussian interior-lifting lemma}

Write $x=(x_1,\ldots,x_n)$ and let
$\Delta_x=\sum_{i=1}^n\partial_{x_i}^2$.

\begin{lemma}[Gaussian lift preserves strict SOS]\label{lem:lift}
If $p\in\cH_{n,2d}$ is a strict sum of squares, then
\[
  P(x,z):=e^{z^2\Delta_x}p(x)
\]
is a strict sum of squares in $\cH_{n+1,2d}$.
\end{lemma}

\begin{proof}
Let $v(x)$ be the monomial vector of $\cH_{n,d}$.  Strict SOS membership
gives
\[
  p(x)=v(x)^{\mathsf T}Gv(x),\qquad G\succ0.
\]
Factoring $G$ shows that
\[
  p=\sum_{j=1}^{N}q_j^2,
  \qquad N=\dim\cH_{n,d},
\]
where $q_1,\ldots,q_N$ form a basis of $\cH_{n,d}$.

Let $\gamma$ be standard Gaussian measure on $\R^n$.  The Gaussian moment
identity gives
\begin{align*}
 P(x,z)
 &=\int_{\R^n}p(x+\sqrt2z\xi)\,d\gamma(\xi)\\
 &=\int_{\R^n}\sum_{j=1}^{N}
       q_j(x+\sqrt2z\xi)^2\,d\gamma(\xi).
 \tag{1}
\end{align*}
Every integrand in (1) is the square of a form in $\cH_{n+1,d}$.  Expanding
its coefficient vector and integrating its rank-one Gram matrix therefore
produces a positive-semidefinite Gram matrix $A$ for $P$.

It remains to show $A\succ0$.  The coefficient vectors of
\[
  q_j(x+z\xi),\qquad 1\leq j\leq N,quad \xi\in\R^n,
  \tag{2}
\]
span $\cH_{n+1,d}$.  Since the $q_j$ form a basis, it suffices to use all
monomials of degree $d$ in place of the $q_j$.  Given a monomial
$x^\alpha z^k$ with $|\alpha|+k=d$, choose a multi-index $\beta$ with
$|\beta|=k$.  Differentiating the translate of $x^{\alpha+\beta}$ gives
\[
 \left.\partial_\xi^\beta
 (x+z\xi)^{\alpha+\beta}\right|_{\xi=0}
 =\frac{(\alpha+\beta)!}{\alpha!}\,x^\alpha z^k.
 \tag{3}
\]
The coefficients of a polynomial map lie in the span of its values, so (3)
proves the spanning claim.  Thus a vector in the kernel of $A$ would be
orthogonal to a spanning set and must vanish.  Hence $A\succ0$, proving that
$P$ is a strict SOS.
\end{proof}

\section{Eventual SOS theorem}

\begin{theorem}\label{thm:general}
Let $f\in\R[x_1,\ldots,x_n]$ have degree $2d$, and let $p=f_{2d}$ be its
highest homogeneous part.  If $p$ is a strict sum of squares, then there is
$T>0$ such that
\[
  e^{t\Delta_x}f\in\SOS(\R^n)
  \qquad\text{for every }t\geq T.
\]
No nonnegativity assumption on $f$ is required.
\end{theorem}

\begin{proof}
Decompose $f=\sum_{k=0}^{2d}f_k$ into homogeneous parts.  For $t>0$ put
$s=t^{-1/2}$.  Homogeneity and the finite heat-series expansion give
\[
 t^{-d}\bigl(e^{t\Delta_x}f\bigr)(\sqrt t\,x)
 =\sum_{k=0}^{2d}s^{2d-k}e^{\Delta_x}f_k(x).
 \tag{4}
\]
The degree-$2d$ homogenization of the right-hand side is
\[
 Q_s(x,z)=\sum_{k=0}^{2d}s^{2d-k}z^{2d-k}
                 e^{z^2\Delta_x}f_k(x).
 \tag{5}
\]
Consequently, as $s\downarrow0$,
\[
  Q_s\longrightarrow Q_0=e^{z^2\Delta_x}p
  \quad\text{coefficientwise in }\cH_{n+1,2d}.
 \tag{6}
\]
Lemma~\ref{lem:lift} places $Q_0$ in the interior of the finite-dimensional
closed cone $\Sigma_{n+1,2d}$.  This interior is open, so (6) implies that
$Q_s\in\Sigma_{n+1,2d}$ for every sufficiently small positive $s$.
Dehomogenizing (5) at $z=1$ and using (4) shows that the normalized polynomial
in (4) is SOS.  Positive scalar multiplication and the invertible scaling
$x\mapsto\sqrt t\,x$ preserve SOS membership.  Thus
$e^{t\Delta_x}f$ is SOS for every sufficiently large $t$.
\end{proof}

\begin{corollary}[Coercive bivariate case]\label{cor:binary}
Let $f\in\R[x,y]$ have degree $2d$ and leading form $p$.  If
\[
  p(u)>0\qquad(u\in\R^2\setminus\{0\}),
\]
then $e^{t\Delta}f$ is SOS for all sufficiently large $t$.  In particular,
Open Problem 5.1 has an affirmative answer whenever the nonnegative input
polynomial has a positive-definite leading form.
\end{corollary}

\begin{proof}
Positive-definite binary forms are precisely the interior points of the cone
of nonnegative binary forms.  Hilbert's binary-form theorem identifies that
cone with $\Sigma_{2,2d}$ \cite{Hilbert}.  Hence $p$ is a strict SOS, and
Theorem~\ref{thm:general} applies.
\end{proof}

\section{Limitations and boundary obstruction}

Corollary~\ref{cor:binary} does not settle the full source problem.  If the
leading binary form $p$ has a real projective zero, then $p$ lies on the
boundary of the binary SOS cone, and $e^{z^2\Delta}p$ lies on a proper face of
the ternary SOS cone.  Coefficient convergence to a boundary point gives no
neighborhood of SOS forms.  Moreover, lower homogeneous pieces can decrease
the multiplicity at an infinity zero, so the normalized path can move through
nested faces rather than remain in a single relative interior.  This is the
specific obstruction left by the proof.

The result nevertheless removes the generic case: among nonnegative binary
leading forms of fixed degree, positive-definite forms constitute the interior
of the cone.  Only the locus with a real projective zero remains.

\section{Novelty check and verification}

The source problem and notation were checked against arXiv:2211.04416, Open
Problem 5.1 on PDF page 26.  A bounded search through 13 August 2026 covered
the local run indexes, the parsed arXiv corpus, exact problem wording, the
source title and arXiv id, and combinations of ``heat'', ``eventual sum of
squares'', ``coercive polynomial'', and ``leading homogeneous form''.  The
later arXiv:2506.16321 cites the source but does not answer this heat-flow
problem.  No explicit prior statement of Theorem~\ref{thm:general} or
Corollary~\ref{cor:binary} was located.  This is bounded evidence, not a claim
of exhaustive novelty.

The proof is exact and has no computational dependency.  Verification checked
the powers of $t$ in (4), the homogenization exponents in (5), the Gaussian
variance in (1), and the monomial-spanning identity (3).  Human review should
focus on those four points and on the equivalence between strict SOS
membership and a positive-definite Gram matrix.

\begin{thebibliography}{9}

\bibitem{CDKM}
R.~E. Curto, P.~J. di Dio, M.~Korda, and V.~Magron,
\emph{Time-dependent moments from partial differential equations and the
time-dependent set of atoms}, arXiv:2211.04416 (2022; revised 2023).

\bibitem{Hilbert}
D.~Hilbert,
\emph{Ueber die Darstellung definiter Formen als Summe von Formenquadraten},
Mathematische Annalen \textbf{32} (1888), 342--350.

\end{thebibliography}

\end{document}

